\subsection{次の接着計算に使う二つの局所障害}
\label{ng15:local-obstructions}

二価 core $i$ の二つの高近傍を $j,k$ とする．$q_{jk}=1$ なら，
core $i$ の二高分母標準形は必ず正の枝である．従って
\[
 S>r,\qquad rT>(r+1)S,\qquad L=T-S+1\geq3.
\]
この結論に対向 $D$ の相殺は不要であり，core $j,k$ の高分母次数を制限しない．
負の枝 $S=T=w<r$ を用いるときは，その core を含む高分母三角形が必要である．
一般の三高分母中心に二高分母の真約数補題を適用する代わりに，
中心側では signed divisors と高辺の合同式を直接比較する．
負の枝の $H=r-w$，$c=s-x$，$e=z-y$ に対して
\[
 Ece=e-c+1,\qquad E=1-\frac{w^2}{H(H+1)},
 \qquad -1<E<\tfrac12,\quad E\ne0
\]
となり，中心間が低辺なら reverse 和が $E-1$ となることが主な障害である．
例外 $e=1$ も node 数 $2+rw/H>6$ により排除される．

全低分母 core $*$ がある場合には，残る四頂点を $I$ とし
\[
 x_{ij}=p^i_{j*},\qquad r_{ij}=p^i_{*j}
\]
と置く．四辺 $i*$ の $D$ 相殺，局所の行和，および高辺の整数 clutching から
\[
 \sum_{i\ne j\in I}x_{ij}=\sum_{i\ne j\in I}r_{ij}=-4,
 \qquad
 \sum_{\substack{i<j\in I\\q_{ij}=1}}
 (x_{ij}+x_{ji}-r_{ij}-r_{ji})=0
\]
を得る．全低 core の自由な entries を固定しなくても使える総和条件である．
diamond では和に低辺が一つだけ残り，二価 core の低低 entries が
反対向き・同じ $L\in\{1,3,4\}$ にそろう．三つの場合は全て排除済みである．
一方，残る $K_4$ 型ではこの和は空となるので，そのまま新しい矛盾にはならない．

\subsection{一高分母の自由度と，二葉三角形の残る割当}
\label{ng15:leaf-data}

一高分母 core は，高 ray $A$ と低 rays $B,C,D$ の適切な標識により
\[
\begin{gathered}
 m=th(h+1),\quad q=m+1,\quad u=th,\quad v=t(h+1),
 \quad h,t\geq1,\\
 P=\begin{pmatrix}
 0&-m&-h-1&h\\
 -1&0&0&0\\
 -u&u-1-c&0&c\\
 v&-v-1-d&d&0
 \end{pmatrix},\qquad c,d\in\ZZ,\quad cd=0
\end{gathered}
\]
と書け，全ての整数 choices に正の weights が存在する．係数は
\[
 (D_A,D_B,D_C,D_D)=(-m-t,\ m+t+1+c+d,\ h-d,\ -h-1-c).
\]
$q=3$ では二つの負の reverse entries が一致し得る．
高 row だけを固定して低低 entries を一つの模型に置き換えず，
この $c,d$ の自由度と必要な低 ray の交換を保持する．

二高分母の低係数は
\[
 D_C=y-x+T-S,\qquad D_D=s-z-1
\]
であり，同時に負とはならない．一高分母の負 reverse を持つ ray $E$ は $B$ または $C$ である．
別の低 ray $J$ から $p_{JE}\leq-1$ を要求される場合，
$E=C$ なら $D_E>0$，残る低 ray から $E$ への entry は零となる．
$E=B$ かつ $D_E\leq0$ なら，残る低 ray からの entry は
$t(h+1)^2-D_E\geq4$ である．
係数相殺と node 条件を組み合わせるときに使える二分である．

二葉三角形を $H=\{01,02,12,03,14\}$ と標識し，二価 core $2$ の高 rays を
$A=0,B=1$ とすると，$C=3,D=4$ は分母上限なしに排除済みである．
交差割当 $C=4,D=3$ でも負の枝は排除された．$S=T=w\ge5$ なら，
node 区間と高辺の clutching により両葉の他葉に向かう係数が正となり，相殺できない．
$w\le4$ の全七模型については，次の有限復元を用いる．
葉 $i$ の高 ray を $H_i$，$R_i=p^2_{iH_i}$，$E_i=D^2_i$ とすると，
ray $2=C$ では $u=th\in[R_i-1,R_i+4]$，$d=h+E_i$，
ray $2=D$ では $v=t(h+1)\in[-4-R_i,1-R_i]$，$c=E_i-h-1$ である．
正約数と，残る自由 entry の六整数の node 区間から全ての葉が有限に復元される．
七模型の葉の全四百八十五組は低辺 $34$ で棄却される．
元の行和・column・二十一の node 配分を用いた独立復元でも，全十四候補集合が一致した．
中心間の分母 $q_{01}$ を列挙する必要はない．

従って残る条件は
\[
 (A,B,C,D)=(0,1,4,3),\qquad S>r,\quad rT>(r+1)S,
 \qquad r\le4\ \text{または}\ z=1.
\]
一般の正の二高分母模型では $x\ge2$ であることも使える．
$r\le4$ でも他の補助整数に上限を置いておらず，$z=1$ の枝も残る．
上の完全葉復元を正の枝へ拡張する場合には，ray $2=B$ が可能となる場合を落とさない．
この support の全排除と存在は未解決である．

\subsection{部分的な低辺条件だけでは完成しない二つの模型}
\label{ng15:partial-patches}

次の例は，部分的な必要条件の非空性を示す資料である．大域 atlas の存在例ではない．
$K_{2,3}$ の三つの二価 core を cyclic に標識し，canonical 高 ray $A$ が
一方の中心を向く場合を $+$，他方を向く場合を $-$ とする．
符号 $(+,+,-)$ に対し，canonical データを
\[
 (r,s,z,S,T)=(4,4,5,2,2),\quad(5,10,2,3,3),\quad(5,10,2,3,3)
\]
と取る．三低辺の node triples は $(0,3,2),(2,2,1),(0,2,3)$ だが，
同じ順の対向低係数和は $7,-10,5$ であり，相殺しない．
中心二 core も与えていない．本版の $K_{2,3}$ 排除は，中心間の低辺と高辺の clutching を加えて成立する．
具体的な全行列・係数・node の再計算は，保存した v14 の node 検算器に含まれる．

diamond の二価 cores $2,3$ と全低 core $4$ についても，
三つの相互低辺の node 条件と $D$ 相殺を全て満たす部分模型がある．
行列の添字順をそれぞれ $(0,1,3,4),(0,1,2,4),(0,1,2,3)$ として
\[
 P_2=\begin{pmatrix}0&-6&1&-2\\-8&0&-2&3\\6&-3&0&-4\\-3&2&0&0\end{pmatrix},\quad
 P_3=\begin{pmatrix}0&-14&-2&5\\-7&0&2&-4\\-5&4&0&0\\2&-2&-1&0\end{pmatrix},
\]
\[
 P_4=\begin{pmatrix}0&0&-1&0\\2&0&0&-3\\0&-1&0&0\\-3&0&2&0\end{pmatrix},
 \quad q_2=(7,7,1,1),\quad q_3=(11,9,1,1),\quad q_4=(1,1,1,1).
\]
係数は順に
\[
 (-2,0,0,2),\qquad(-1,3,0,-2),\qquad(0,0,-2,2)
\]
であり，低辺 $23,24,34$ の node triples はそれぞれ
$(0,0,5)$，$(4,0,1)$，$(2,3,0)$ である．
全てに正の weights がある．しかし全低 hub の残差
\[
 p^2_{34}+p^3_{24}-p^2_{43}-p^3_{42}=-3
\]
は零でないので，二つの中心 core を加えた大域完成は不可能である．
行和・pair products・正 weights・係数・nodes の厳密検算は
\nolinkurl{diamond_obstruction_verification_v15.json} に保存した．
部分模型の存在を保ちながら，追加する大域条件のどこで矛盾するかを区別するための例である．
